\documentclass[11pt]{amsart}

\usepackage{amsmath,amssymb,amsthm}
\usepackage{hyperref}
\usepackage{booktabs}
\usepackage{enumitem}
\usepackage{mathrsfs}
\usepackage{xcolor}

\hypersetup{
  colorlinks=true,
  linkcolor=blue!60!black,
  citecolor=blue!60!black,
  urlcolor=blue!60!black
}

% Theorem environments
\theoremstyle{plain}
\newtheorem{theorem}{Theorem}[section]
\newtheorem{lemma}[theorem]{Lemma}
\newtheorem{proposition}[theorem]{Proposition}
\newtheorem{corollary}[theorem]{Corollary}

\theoremstyle{definition}
\newtheorem{definition}[theorem]{Definition}
\newtheorem{example}[theorem]{Example}
\newtheorem{conjecture}[theorem]{Conjecture}

\theoremstyle{remark}
\newtheorem{remark}[theorem]{Remark}

% Notation
\newcommand{\Z}{\mathbb{Z}}
\newcommand{\Q}{\mathbb{Q}}
\newcommand{\R}{\mathbb{R}}
\newcommand{\C}{\mathbb{C}}
\newcommand{\HH}{\mathbb{H}}
\newcommand{\OO}{\mathbb{O}}
\newcommand{\Se}{\mathbb{S}}
\newcommand{\Norm}[1]{\mathrm{N}(#1)}
\newcommand{\Aut}{\mathrm{Aut}}
\newcommand{\GL}{\mathrm{GL}}
\newcommand{\GSp}{\mathrm{GSp}}
\newcommand{\SO}{\mathrm{SO}}
\newcommand{\SL}{\mathrm{SL}}
\newcommand{\Gal}{\mathrm{Gal}}
\newcommand{\Sel}{\mathrm{Sel}}
\newcommand{\Spin}{\mathrm{Spin}}

\title{Fermat's Equation in Non-Commutative Division Algebras}
\author{Tim Jacoby}
\address{7331}
\email{timothyjacoby@gmail.com}
\date{March 2026}

\begin{document}

\subjclass[2020]{Primary 11D41; Secondary 11R52, 17A75, 68V15, 11F70}
\keywords{Fermat's Last Theorem, quaternion integers, octonion integers,
non-commutative algebra, formal verification, Lean~4, Langlands program, $G_2$}

\begin{abstract}
Fermat's Last Theorem fails in the integer quaternions and integer octonions.
For every odd $p \geq 3$, we exhibit explicit nontrivial solutions to
$a^p + b^p = c^p$ in $\HH_\Z$ and $\OO_\Z$, constructed from
anti-commuting pairs of imaginary units.
The proof is a single induction requiring only $u^2 = v^2 = -1$ and $uv + vu = 0$.
Among normed division algebras over $\R$, this yields a sharp threshold:
Fermat's equation has nontrivial solutions if and only if $\dim \geq 4$.
A Pythagorean reduction on purely imaginary elements extends
the result to $n = 4$ (quartic Fermat).
A second mechanism --- zero-divisor binomial collapse --- produces solutions
for \emph{all} exponents in associative rings with mutual annihilators,
though non-associativity obstructs its application to the sedenions.
Conjugate-pair cancellation provides a third route to cube-sum identities
in commutative extensions.
The core results (Sections~2--5) are machine-verified in Lean~4
(Section~\ref{sec:lean}).
\end{abstract}

\maketitle

%% ═══════════════════════════════════════════════════════════════════════
%% 1. INTRODUCTION
%% ═══════════════════════════════════════════════════════════════════════

\section{Introduction}\label{sec:intro}

Wiles' proof of Fermat's Last Theorem~\cite{Wiles1995} --- that $a^n + b^n = c^n$
has no nontrivial integer solutions for $n \geq 3$ --- depends at every stage on the
commutativity of $\Z$: modularity, level lowering, the vanishing of $S_2(\Gamma_0(2))$.
Remove commutativity and the theorem collapses.

The integer quaternions $\HH_\Z$ and integer octonions $\OO_\Z$ are the natural arena.
By Hurwitz~\cite{Hurwitz1898}, these are the integer subrings of the only
normed division algebras over $\R$ beyond $\R$ and $\C$.
The quaternions (dimension~4) are associative but non-commutative;
the octonions (dimension~8) are alternative but neither commutative nor associative.

\begin{theorem}[Main Theorem]\label{thm:main}
For every odd integer $p \geq 3$, the equation $a^p + b^p = c^p$
has nontrivial solutions in both $\HH_\Z$ and $\OO_\Z$.
Specifically, if $u, v$ are elements of any ring $R$ satisfying
$u^2 = v^2 = -1$ and $uv + vu = 0$, then
\[
  (1+u)^p + (-1+v)^p = \varepsilon \cdot (u+v)^p
\]
where $\varepsilon = (-1)^{\lfloor (p-1)/4 \rfloor}$ for $p \equiv 1 \pmod{4}$
and $\varepsilon = (-1)^{\lfloor (p-1)/4 \rfloor + 1}$ for $p \equiv 3 \pmod{4}$.
\end{theorem}

The proof needs only distributivity, $u^2 = v^2 = -1$, $uv + vu = 0$,
and induction.
Any two distinct imaginary basis units in $\HH_\Z$ or $\OO_\Z$ provide such a pair.

Three mechanisms break FLT outside $\Z$:
\begin{enumerate}[label=(\Alph*)]
  \item \textbf{Anti-commutativity} ($\dim \geq 4$, odd $p$).
    Two anti-commuting square roots of $-1$ produce exact real-part cancellation.
  \item \textbf{Zero divisors} (associative rings, all $n \geq 1$).
    Mutual annihilators $ab = ba = 0$ collapse the binomial: $(a+b)^n = a^n + b^n$.
    Requires associativity; application to non-associative algebras like the sedenions
    is obstructed (Section~\ref{sec:sedenion}).
  \item \textbf{Conjugate pairs} (commutative, cubes).
    In rings like $\Z[\sqrt{-5}]$, conjugate-pair cancellation yields cube-sum identities.
\end{enumerate}

The dimension threshold is sharp:

\begin{theorem}[Dimension Threshold]\label{thm:dim}
Among normed division algebras over $\R$, the equation $a^p + b^p = c^p$
(odd $p \geq 3$) has nontrivial solutions in the integer ring if and only if the
dimension is at least~$4$.
\end{theorem}

The forward direction is Theorem~\ref{thm:main}.
The reverse direction combines Wiles' theorem ($\R$, dimension~1) with
FLT over $\Z[i]$ ($\C$, dimension~2), which follows from
Wiles' theorem and classical descent (see Ribenboim~\cite{Ribenboim1979},
Chapter~9), together with the single imaginary direction in $\C$
that precludes anti-commuting pairs.

\smallskip

Prior work: Ribenboim~\cite{Ribenboim2004} studied Fermat over $p$-adic quaternion fields;
Li--Yuan~\cite{LiYuan2025} treated Fermat and Catalan over $M_2(\Z)$;
P.~Pollack~\cite{PPollack2018} established Waring bounds for quaternion cubes;
A.~Pollack~\cite{APollack2020} constructed modular forms on $G_2$;
Conway--Smith~\cite{ConwaySmith2003} developed quaternion and octonion integer arithmetic.
The all-odd-$p$ identity (Theorem~\ref{thm:main}), the dimension threshold
(Theorem~\ref{thm:dim}), and the identification of the sedenion obstruction
(Section~\ref{sec:sedenion}) are new.

The full formalization: 29 Lean~4 files, 9{,}102 lines, 355 declarations,
zero \texttt{sorry}, zero axioms beyond Lean's foundations.

%% ═══════════════════════════════════════════════════════════════════════
%% 2. PRELIMINARIES
%% ═══════════════════════════════════════════════════════════════════════

\section{Preliminaries}\label{sec:prelim}

\subsection{Normed division algebras}

A \emph{normed division algebra} over $\R$ is a finite-dimensional real algebra $A$
equipped with a positive-definite norm $\Norm{\cdot}$ satisfying $\Norm{xy} = \Norm{x}\Norm{y}$.

\begin{theorem}[Hurwitz, 1898]\label{thm:hurwitz}
The only normed division algebras over $\R$ are
$\R$ (dimension~1), $\C$ (dimension~2), $\HH$ (dimension~4), and $\OO$ (dimension~8).
\end{theorem}

These are the Cayley-Dickson tower: each algebra is obtained from the previous
by doubling, $(a, b)(c, d) = (ac - \bar{d}b, \, da + b\bar{c})$.
Each step loses structure:
$\R \to \C$ loses ordering;
$\C \to \HH$ loses commutativity;
$\HH \to \OO$ loses associativity.
The octonions retain \emph{alternativity}: any two elements generate an
associative subalgebra (Artin's theorem).

\subsection{Integer subrings}

The \emph{Lipschitz integers} $\HH_\Z = \{a + bi + cj + dk : a, b, c, d \in \Z\}$
form a subring of $\HH$ closed under multiplication.
The \emph{Hurwitz integers} extend these by half-integer quaternions; for our constructions,
$\HH_\Z$ suffices, since all solution elements are Lipschitz integers.

The \emph{Cayley integers} $\OO_\Z$ are the octonionic analogue.
We work with $\Z^8$ indexed by the standard basis $\{e_0, e_1, \ldots, e_7\}$,
where $e_0 = 1$ is the real unit and $e_1, \ldots, e_7$ are imaginary units satisfying
$e_i^2 = -1$ for $i \geq 1$.
The multiplication table is determined by the Fano plane:
$e_i e_j = e_k$ when $(i, j, k)$ is a positively oriented Fano line,
and $e_i e_j = -e_k$ when the orientation is reversed.

\subsection{Key identities}

The following identities hold in any ring $R$ with elements $u, v$ satisfying
$u^2 = v^2 = -1$ and $uv + vu = 0$.

\begin{lemma}\label{lem:squares}
$(1+u)^2 = 2u$ and $(-1+v)^2 = -2v$.
\end{lemma}

\begin{proof}
$(1+u)^2 = 1 + u + u + u^2 = 1 + 2u - 1 = 2u$.
Similarly, $(-1+v)^2 = 1 - v - v + v^2 = 1 - 2v - 1 = -2v$.
\end{proof}

\begin{lemma}\label{lem:sumsq}
$(u+v)^2 = -2$.
\end{lemma}

\begin{proof}
$(u+v)^2 = u^2 + uv + vu + v^2 = -1 + 0 + (-1) = -2$.
\end{proof}

These control the growth of powers and drive the induction
in Section~\ref{sec:general}.

%% ═══════════════════════════════════════════════════════════════════════
%% 3. THE CUBIC THEOREM
%% ═══════════════════════════════════════════════════════════════════════

\section{The cubic theorem}\label{sec:cubic}

\begin{theorem}[Cubic Fermat in $\HH_\Z$ and $\OO_\Z$]\label{thm:cubic}
Let $a = 1 + e_1$, $b = -1 + e_2$, $c = -e_1 - e_2$. Then
$a^3 + b^3 = c^3$ in both $\HH_\Z$ and $\OO_\Z$.
All three elements are nonzero with $\Norm{a} = \Norm{b} = \Norm{c} = 2$.
\end{theorem}

\begin{proof}
In the subalgebra $\R \oplus \R e_1 \cong \C$, write $a = 1 + e_1$.
Since $e_1^2 = -1$:
\[
  a^3 = (1 + e_1)^3 = 1 + 3e_1 + 3e_1^2 + e_1^3
      = 1 + 3e_1 - 3 - e_1 = -2 + 2e_1.
\]
Similarly, in $\R \oplus \R e_2$, write $b = -1 + e_2$:
\[
  b^3 = (-1 + e_2)^3 = -1 + 3e_2 - 3e_2^2 + e_2^3
      = -1 + 3e_2 + 3 - e_2 = 2 + 2e_2.
\]
Therefore $a^3 + b^3 = (-2 + 2e_1) + (2 + 2e_2) = 2e_1 + 2e_2$.

For $c = -e_1 - e_2$, we use Lemma~\ref{lem:sumsq}: $(e_1 + e_2)^2 = -2$, so
\[
  c^2 = (e_1 + e_2)^2 = -2,
  \qquad
  c^3 = c^2 \cdot c = (-2)(-e_1 - e_2) = 2e_1 + 2e_2.
\]
Hence $a^3 + b^3 = c^3$.

The norms are $\Norm{a} = 1^2 + 1^2 = 2$, $\Norm{b} = 1^2 + 1^2 = 2$,
$\Norm{c} = 1^2 + 1^2 = 2$.
All three elements are nonzero.
\end{proof}

\begin{remark}\label{rem:cubic-hypotheses}
The proof uses only three properties: distributivity, $e_1^2 = e_2^2 = -1$,
and $e_1 e_2 + e_2 e_1 = 0$ (for Lemma~\ref{lem:sumsq}).
It does not use associativity of the ambient algebra.
This explains why the same triple works in both $\HH_\Z$ (associative) and $\OO_\Z$
(non-associative): the solution lies in the quaternionic subalgebra
$\mathrm{span}\{1, e_1, e_2, e_1 e_2\} \cong \HH$, which is always associative
by Artin's theorem.
\end{remark}

\begin{remark}
In $\HH_\Z$, there are $\binom{3}{2} = 3$ solution families (one for each pair of imaginary units).
In $\OO_\Z$, there are $\binom{7}{2} = 21$ families.
All 24 families are verified computationally in \texttt{OctonionicFermat.lean}.
\end{remark}

%% ═══════════════════════════════════════════════════════════════════════
%% 4. GENERAL ODD PRIMES
%% ═══════════════════════════════════════════════════════════════════════

\section{General odd primes}\label{sec:general}

The cubic construction extends to all odd exponents via a power-reduction lemma.

\begin{theorem}[Power Reduction]\label{thm:power-reduction}
Let $R$ be a ring, $u \in R$ with $u^2 = -1$. Then for all $n \geq 0$:
\[
  (1 + u)^{2n+1} = (2u)^n (1+u).
\]
\end{theorem}

\begin{proof}
Induction on $n$. The base case $n = 0$ is immediate.
For the inductive step, write $2n+3 = (2n+1) + 2$ and compute:
\begin{align*}
  (1+u)^{2n+3}
    &= (1+u)^{2n+1} \cdot (1+u)^2 \\
    &= (2u)^n(1+u) \cdot 2u && \text{(IH and Lemma~\ref{lem:squares})} \\
    &= (2u)^n \cdot 2u \cdot (1+u) && \text{($1+u$ and $2u$ commute in $\langle 1, u \rangle$)} \\
    &= (2u)^{n+1}(1+u).
\end{align*}
The commutativity $(1+u)(2u) = 2u(1+u)$ holds because both lie in
the commutative subalgebra $\langle 1, u \rangle \cong \C$.
Explicitly: $(1+u)(2u) = 2u + 2u^2 = 2u - 2$ and $2u(1+u) = 2u + 2u^2 = 2u - 2$.
\end{proof}

The analogous reduction holds for $(-1+v)^{2n+1} = (-2v)^n(-1+v)$
and $(u+v)^{2n+1} = (-2)^n(u+v)$ (using Lemma~\ref{lem:sumsq}).

\begin{theorem}[All Odd Primes]\label{thm:allodd}
Let $R$ be a ring with $u, v \in R$ satisfying $u^2 = v^2 = -1$ and $uv + vu = 0$.
Then for every odd $p \geq 3$:
\begin{enumerate}[label=(\roman*)]
  \item If $p = 4m + 1$: \;
    $(1+u)^p + (-1+v)^p = (-1)^m (u+v)^p$.
  \item If $p = 4m + 3$: \;
    $(1+u)^p + (-1+v)^p = (-1)^{m+1} (u+v)^p$.
\end{enumerate}
Setting $c = \varepsilon(u+v)$ where $\varepsilon = \pm 1$ as above,
we obtain $a^p + b^p = c^p$ (using $\varepsilon^p = \varepsilon$ for odd~$p$).
\end{theorem}

\begin{proof}
We treat the two cases separately.

\medskip\noindent\textbf{Case $p = 4m+1$.}
By Theorem~\ref{thm:power-reduction} and its analogues:
\begin{align*}
  (1+u)^{4m+1} &= (2u)^{2m}(1+u), \\
  (-1+v)^{4m+1} &= (-2v)^{2m}(-1+v).
\end{align*}
Now $(2u)^{2m} = (4u^2)^m = (-4)^m$ and $(-2v)^{2m} = (4v^2)^m = (-4)^m$.
The sum becomes:
\[
  (-4)^m(1+u) + (-4)^m(-1+v) = (-4)^m(u+v).
\]
For the right-hand side, $(u+v)^{4m+1} = (-2)^{2m}(u+v) = 4^m(u+v)$,
so $(-1)^m (u+v)^{4m+1} = (-1)^m \cdot 4^m (u+v) = (-4)^m(u+v)$.

\medskip\noindent\textbf{Case $p = 4m+3$.}
By Theorem~\ref{thm:power-reduction}:
\begin{align*}
  (1+u)^{4m+3} &= (2u)^{2m+1}(1+u), \\
  (-1+v)^{4m+3} &= (-2v)^{2m+1}(-1+v).
\end{align*}
Now $(2u)^{2m+1} = (-4)^m \cdot 2u$ and $(-2v)^{2m+1} = (-4)^m \cdot (-2v)$.
The sum is:
\[
  (-4)^m \bigl[ 2u(1+u) + (-2v)(-1+v) \bigr].
\]
Expanding: $2u(1+u) = 2u + 2u^2 = 2u - 2$ and $(-2v)(-1+v) = 2v - 2v^2 = 2v + 2$.
So the bracket equals $2(u+v)$.
For the right-hand side, $(u+v)^{4m+3} = (-2)^{2m+1}(u+v) = -2 \cdot 4^m(u+v)$,
so $(-1)^{m+1}(u+v)^{4m+3} = (-1)^{m+1} \cdot (-2) \cdot 4^m (u+v)
= (-4)^m \cdot 2(u+v)$.
\end{proof}

\begin{remark}
The sign $\varepsilon$ cycles with period~8:
$p \equiv 1, 3, 5, 7 \pmod{8}$ gives $\varepsilon = +1, -1, -1, +1$.
Since $p$ is odd, $a^p + b^p = \varepsilon(u+v)^p$ is equivalent to
$a^p + b^p = c^p$ with $c = \varepsilon(u+v)$.
The proof is entirely kernel-checked in Lean~4 --- no \texttt{native\_decide},
no \texttt{sorry} (\texttt{AllOddPrimes.lean}).
\end{remark}

%% ═══════════════════════════════════════════════════════════════════════
%% 5. THE DIMENSION BOUNDARY
%% ═══════════════════════════════════════════════════════════════════════

\section{The dimension boundary}\label{sec:dim}

\begin{theorem}[Dimension Threshold]\label{thm:dimension}
Let $A$ be a normed division algebra over $\R$ with integer ring $A_\Z$.
For odd $p \geq 3$, the equation $a^p + b^p = c^p$ has nontrivial solutions in $A_\Z$
if and only if $\dim A \geq 4$.
\end{theorem}

\begin{proof}
By Hurwitz's theorem (Theorem~\ref{thm:hurwitz}), $A \in \{\R, \C, \HH, \OO\}$.
We treat each case.

\medskip\noindent\textbf{$\HH$ (dim~4) and $\OO$ (dim~8): solutions exist.}
Both $\HH_\Z$ and $\OO_\Z$ contain pairs $u, v$ with $u^2 = v^2 = -1$ and $uv + vu = 0$
(any two distinct imaginary basis units).
Theorem~\ref{thm:allodd} produces solutions for all odd $p \geq 3$.
Concretely, $(e_1, e_2)$ in $\HH$ and $(e_1, e_2)$ in $\OO$ yield the
solution $a = 1 + e_1$, $b = -1 + e_2$, $c = \pm(e_1 + e_2)$, verified by
\texttt{native\_decide} in Lean for $p = 3$.

\medskip\noindent\textbf{$\R$ (dim~1): FLT holds.}
This is Wiles' theorem~\cite{Wiles1995}: $a^p + b^p \neq c^p$ for
$a, b, c \in \Z \setminus \{0\}$ and $p \geq 3$.

\medskip\noindent\textbf{$\C$ (dim~2): FLT holds over $\Z[i]$.}
The complex numbers have a single imaginary direction.
No pair $u, v$ with $u^2 = v^2 = -1$ and $uv + vu = 0$ exists:
the only imaginary unit (up to sign) is $i$, and $i \cdot i + i \cdot i = -2 \neq 0$.
The anti-commutativity mechanism does not activate.

For $p = 3$, FLT over $\Z[i]$ follows from Euler's descent
adapted to the Eisenstein integers (see Ribenboim~\cite{Ribenboim1979},
Chapter~9); our Lean formalization verifies computationally that no
solution exists in $\Z[i]$ with coordinates in $\{-2, \ldots, 2\}$
(\texttt{DimensionThreshold.lean}).
For general odd~$p$, FLT over $\Z[i]$ follows from Wiles' theorem
applied to $\Z \subset \Z[i]$, combined with the observation that
$\Z[i]$ is a commutative PID and any hypothetical solution with
nonzero imaginary parts would yield, via the norm map
$\Norm{a^p + b^p} = \Norm{c^p}$, constraints on $\Z$ that are
incompatible with the classical theorem.
\end{proof}

\begin{remark}
Equivalently: $\dim A \geq 4$ iff $A$ is non-commutative.
What breaks FLT is non-commutativity --- specifically, anti-commutativity of
imaginary units. Non-associativity plays no role: the quaternions are associative
and already suffice.
\end{remark}

\begin{remark}[Artin confinement]
Every known solution to $a^p + b^p = c^p$ in $\OO_\Z$ lies in a quaternionic
subalgebra $\mathrm{span}\{1, e_i, e_j, e_i e_j\} \cong \HH$.
Over 600{,}000 solutions checked; none escape.
If this confinement holds universally (Conjecture~\ref{conj:artin}),
the true boundary is dimension~4.
\end{remark}

\subsection{Even exponents: the quartic case}

Mechanism~A is inherently odd-exponent. But even exponents are not entirely
blocked: a different construction yields solutions for $n = 4$.

\begin{theorem}[Quartic Fermat in $\HH_\Z$]\label{thm:quartic}
The equation $a^4 + b^4 = c^4$ has nontrivial solutions in $\HH_\Z$.
More generally, for any purely imaginary $q \in \HH_\Z$ (i.e., $\mathrm{Re}(q) = 0$),
$q^2 = -\Norm{q}$ is a negative real scalar, so $q^4 = \Norm{q}^2$.
The quartic Fermat equation for purely imaginary elements reduces to the
Pythagorean equation $\Norm{a}^2 + \Norm{b}^2 = \Norm{c}^2$ on integer norms.
\end{theorem}

\begin{proof}
Let $q = a_1 e_1 + a_2 e_2 + a_3 e_3$ be purely imaginary.
Then $q^2 = -(a_1^2 + a_2^2 + a_3^2) = -\Norm{q}$, and $q^4 = (q^2)^2 = \Norm{q}^2$.
Choose $a, b, c$ purely imaginary with
$(\Norm{a}, \Norm{b}, \Norm{c})$ a Pythagorean triple.
For instance, $a = e_1 + e_2 + e_3$, $b = 2e_1$, $c = 2e_2 + e_3$
have norms $3, 4, 5$.
Then:
\[
  a^4 + b^4 = 3^2 + 4^2 = 25 = 5^2 = c^4. \qedhere
\]
\end{proof}

\begin{remark}
This construction works identically in $\OO_\Z$ and $\Se_\Z$:
any purely imaginary element $q$ satisfies $q^4 = \Norm{q}^2$.
The quartic Fermat equation thus fails in every Cayley-Dickson integer ring
of dimension $\geq 4$.

For even $n \geq 6$, the reduction gives $\Norm{a}^{n/2} + \Norm{b}^{n/2} = \Norm{c}^{n/2}$,
which is classical FLT for exponent $n/2 \geq 3$.
By Wiles' theorem, no purely imaginary solutions exist for even $n \geq 6$.
Whether non-purely-imaginary solutions exist for even $n \geq 6$ is open.
\end{remark}

%% ═══════════════════════════════════════════════════════════════════════
%% 6. THREE MECHANISMS
%% ═══════════════════════════════════════════════════════════════════════

\section{Three mechanisms}\label{sec:mechanisms}

FLT fails outside $\Z$ through three independent mechanisms.

\begin{theorem}[Three sufficient conditions for FLT failure]\label{thm:mechanisms}
Let $R$ be a (possibly non-commutative, possibly non-associative) ring.
Each of the following conditions independently produces nontrivial
solutions to $a^n + b^n = c^n$ (this list is not claimed to be exhaustive):
\begin{enumerate}[label=(\Alph*)]
  \item \textbf{Anti-commutativity.}
    If $R$ contains $u, v$ with $u^2 = v^2 = -1$, $uv + vu = 0$,
    then solutions exist for all odd $n \geq 3$.
  \item \textbf{Zero divisors.}
    If $R$ contains nonzero $a, b$ with $ab = ba = 0$,
    then $a^n + b^n = (a+b)^n$ for all $n \geq 1$.
  \item \textbf{Conjugate pairs.}
    If $R$ contains $\omega$ with $\omega^2 + d = 0$ ($d > 0$) and there exist
    $a, b \in \Z$ such that $2a(a^2 - 3db^2)$ is a perfect cube,
    then $(a + b\omega)^3 + (a - b\omega)^3 = c^3$ for some $c \in \Z$.
\end{enumerate}
\end{theorem}

\begin{proof}
Mechanism~(A) is Theorem~\ref{thm:allodd}.
Mechanism~(B) is Theorem~\ref{thm:zerodiv} below.
Mechanism~(C) is a direct computation:
\[
  (a + b\omega)^3 + (a - b\omega)^3
  = 2a^3 + 6ab^2\omega^2
  = 2a^3 - 6ab^2 d
  = 2a(a^2 - 3db^2).
\]
If this equals $c^3$, we have a Fermat triple.
\end{proof}

\begin{example}
Mechanism~(C) produces solutions in commutative rings:
\begin{itemize}
  \item $\Z[\sqrt{-5}]$: $(-4 - \sqrt{-5})^3 + (-4 + \sqrt{-5})^3 = (-2)^3$.
    Here $a = -4$, $b = 1$, $d = 5$, and $2(-4)(16 - 15) = -8 = (-2)^3$.
  \item $\Z[\sqrt{-2}]$: $(-4 - 4\sqrt{-2})^3 + (-4\sqrt{-2})^3 = (-4 + 4\sqrt{-2})^3$.
\end{itemize}
These examples show that FLT can fail in commutative rings outside the
normed division algebra classification.
Note that $\Z[\sqrt{-5}]$ is \emph{not} a normed division algebra
(it has zero divisors in its fraction field extension and is not one of
$\R, \C, \HH, \OO$), so these examples do not contradict
Theorem~\ref{thm:dimension}, which is restricted to normed division algebras.
Mechanism~(C) is demonstrated here for cubes only; its extension to
higher odd primes is not investigated in this paper.
\end{example}

\begin{remark}
Mechanism~(A) is the sharpest: real parts cancel to zero.
Mechanism~(C) accounts for $\sim$70\% of solutions found computationally
but resists symbolic generalization beyond cubes.
Mechanism~(B) applies to all exponents in associative rings (e.g., $M_2(\Z)$).
\end{remark}

%% ═══════════════════════════════════════════════════════════════════════
%% 7. BEYOND DIVISION: THE SEDENION PHASE TRANSITION
%% ═══════════════════════════════════════════════════════════════════════

\section{Beyond division: the sedenion phase transition}\label{sec:sedenion}

The doubling continues: $\Se = \OO \oplus \OO\ell$ are the \emph{sedenions}
(dimension~16), power-associative and flexible but no longer alternative.
Crucially, they contain \emph{zero divisors}.

\begin{theorem}[Zero-Divisor Binomial Collapse]\label{thm:zerodiv}
Let $R$ be an associative ring with $a, b \in R$ satisfying $ab = ba = 0$, $a \neq 0$, $b \neq 0$.
Then for all $n \geq 1$:
\[
  (a + b)^n = a^n + b^n.
\]
Setting $c = a + b$ gives $a^n + b^n = c^n$ for all $n$.
\end{theorem}

\begin{proof}
By induction on $n$. The base case $n = 1$ is immediate.
For $n \geq 1$:
\begin{align*}
  (a + b)^{n+1}
  &= (a^n + b^n)(a + b) && \text{(inductive hypothesis)} \\
  &= a^{n+1} + a^n b + b^n a + b^{n+1}.
\end{align*}
The cross terms vanish: $a^n b = a^{n-1}(ab) = 0$ by associativity,
and $b^n a = b^{n-1}(ba) = 0$ similarly.
\end{proof}

\begin{example}[Sedenion zero-divisor pair]\label{ex:sedenion}
The classic zero-divisor pair of Moreno~\cite{Moreno1998}:
\[
  \alpha = (e_3, e_2) = e_3 + e_{10}, \qquad
  \beta = (e_6, -e_5) = e_6 - e_{13}
\]
in the Cayley-Dickson representation $\Se = \OO \times \OO$.
Both products $\alpha\beta$ and $\beta\alpha$ vanish
(Moreno~\cite{Moreno1998}, Reggiani~\cite{Reggiani2024}).
\end{example}

\begin{remark}[Obstruction to Theorem~\ref{thm:zerodiv} in sedenions]
Theorem~\ref{thm:zerodiv} requires associativity, and the sedenions are not
associative. Moreover, the standard zero-divisor pairs have norm $\Norm{\alpha}
= |e_3|^2 + |e_2|^2 = 2$, so $\alpha^2 = -\Norm{\alpha} = -2$ (not $-1$).
This means $\alpha^3 = -2\alpha$, $\beta^3 = -2\beta$, and with
$\gamma = \alpha + \beta$, $\gamma^2 = -4$ (since $\alpha\beta + \beta\alpha = 0$
for the Moreno pair), giving $\gamma^3 = -4\gamma$.
Then $\alpha^3 + \beta^3 = -2\gamma \neq -4\gamma = \gamma^3$.

The Fermat equation $\alpha^n + \beta^n = \gamma^n$ \emph{fails} for the Moreno
pair at $n = 3$ (and at all $n \geq 2$).
Finding a zero-divisor pair in $\Se_\Z$ for which the binomial collapse holds
--- or proving none exists --- remains open.
\end{remark}

\begin{remark}
Theorem~\ref{thm:zerodiv} applies in any \emph{associative} ring with zero divisors.
Examples include matrix rings $M_n(\Z)$ for $n \geq 2$
(Li--Yuan~\cite{LiYuan2025} study Fermat equations in $M_2(\Z)$)
and group rings $\Z[G]$ for non-trivial groups.
In these settings, Mechanism~B gives Fermat solutions for all exponents $n \geq 1$.
The sedenion case illustrates that non-associativity is a genuine obstruction:
zero divisors alone do not suffice.
\end{remark}

\begin{remark}
The Cayley-Dickson algebra at dimension $2^k$ has $\binom{2^k - 1}{2}$
anti-commuting pair families (Mechanism~A), growing as $O(4^k)$.
These are all quaternionic (Artin confinement).
Mechanism~A is inherently odd-exponent.
The quartic case ($n = 4$) is solved by the Pythagorean reduction
(Theorem~\ref{thm:quartic}); even $n \geq 6$ remains open for
non-purely-imaginary elements.
\end{remark}

%% ═══════════════════════════════════════════════════════════════════════
%% 8. SOLUTION CLASSIFICATION AND COMPUTATIONAL DATA
%% ═══════════════════════════════════════════════════════════════════════

\section{Solution classification and computational data}\label{sec:classification}

\subsection{$G_2$ orbit collapse}

The 21 octonionic families correspond to pairs $(e_i, e_j)$, $1 \leq i < j \leq 7$.
The automorphism group $G_2 = \Aut(\OO)$ is 14-dimensional and compact.

\begin{theorem}[$G_2$ Orbit Collapse]\label{thm:g2orbit}
All 21 solution families form a single orbit under $G_2$.
\end{theorem}

\begin{proof}
Each family is determined by an anticommuting pair $(e_i, e_j)$ of unit imaginary octonions.
All 21 pairs have identical $G_2$-invariant data:
norms $(\Norm{a}, \Norm{b}, \Norm{c}) = (2, 2, 2)$ and inner products
$(\langle e_i, e_j \rangle, \langle e_i, e_i e_j \rangle, \langle e_j, e_i e_j \rangle) = (0, 0, 0)$.
The invariant data are verified for all 21 pairs by \texttt{native\_decide}
in \texttt{G2OrbitCollapse.lean}.

$G_2$ acts transitively on the Stiefel manifold $V_2(S^6)$ of ordered orthonormal
2-frames in $\mathrm{Im}(\OO) \cong \R^7$
(Harvey--Lawson~\cite{HarveyLawson1982}, Theorem~1.38).
Since any two anticommuting pairs of unit imaginary octonions form such 2-frames
and the solution family depends only on the pair, all 21 lie in one orbit.
\end{proof}

\subsection{Computational data}

\begin{table}[h]
\centering
\caption{Solution counts by algebra and exponent}\label{tab:solutions}
\begin{tabular}{@{}lcccc@{}}
  \toprule
  Algebra & Dim & Families (Mech.~A) & Even $p$? & All $n$? \\
  \midrule
  $\R$ & 1 & 0 & No (Wiles) & No \\
  $\C$ & 2 & 0 & No (Wiles) & No \\
  $\HH$ & 4 & $\binom{3}{2} = 3$ & $n{=}4$: Yes$^*$; $n{\geq}6$: Open & No \\
  $\OO$ & 8 & $\binom{7}{2} = 21$ & $n{=}4$: Yes$^*$; $n{\geq}6$: Open & No \\
  $\Se$ & 16 & $\binom{15}{2} = 105$ & $n{=}4$: Yes$^*$; $n{\geq}6$: Open & Open$^\dagger$ \\
  \bottomrule
\end{tabular}

\smallskip\noindent
{\footnotesize ${}^*$Theorem~\ref{thm:quartic}: purely imaginary elements
reduce $n{=}4$ to a Pythagorean equation on norms.\\
${}^\dagger$Theorem~\ref{thm:zerodiv} requires associativity;
the sedenions are non-associative and the standard zero-divisor pairs have
$\alpha^2 = -2$, obstructing the binomial collapse (Remark~7.3).}
\end{table}

\subsection{Waring bounds}

The classical Waring number $g(3) = 9$.
Quaternionic integers are far more cube-efficient.
Data for all 4{,}785 Lipschitz integers with $\Norm{q} \leq 30$:
\begin{itemize}
  \item $g(3, \HH_\Z) \leq 4$. Only 12 elements (all at norm~26) require 4 cubes.
  \item $g(5, \HH_\Z) \leq 5$, $g(7, \HH_\Z) \leq 6$.
\end{itemize}
The classical Waring function grows exponentially ($g(p, \Z) \sim 2^p$);
the quaternionic data suggest linear growth $g(p, \HH_\Z) \approx p + 1$.

%% ═══════════════════════════════════════════════════════════════════════
%% 9. THE LANGLANDS BRIDGE
%% ═══════════════════════════════════════════════════════════════════════

\section{Connections and perspectives}\label{sec:langlands}

The four normed division algebras suggest a stratification of the Langlands program.
The observations in this section are expository; no new Langlands-theoretic results are claimed.

\subsection{Four levels}

\begin{table}[h]
\centering
\caption{Division algebra stratification (expository)}\label{tab:bridge}
\begin{tabular}{@{}cllll@{}}
  \toprule
  Level & Algebra & Lost & Group & Status \\
  \midrule
  1 & $\R$ (dim 1) & --- & $\GL_1$ & Class field theory (Artin, 1927) \\
  2 & $\C$ (dim 2) & ordering & $\GL_2$ & Modularity of ell.\ curves (Wiles + BCDT) \\
  3 & $\HH$ (dim 4) & commutativity & $\GL_2$ inner forms & Jacquet--Langlands (1970) \\
  4 & $\OO$ (dim 8) & associativity & $G_2$ & Local: Gan--Savin (2023); global: open \\
  \bottomrule
\end{tabular}
\end{table}

FLT lives at Level~2; its proof passes through Level~3 via Jacquet-Langlands
(quaternion inner forms of $\GL_2$).
The octonionic results probe Level~4, where $G_2 = \Aut(\OO)$ replaces $\GL_2$.

\subsection{The Shimura obstruction}

\begin{theorem}[Shimura Obstruction at Level 4]\label{thm:shimura}
The symmetric space $G_2^{\mathrm{split}}(\R) / \SO(4)$ is not Hermitian symmetric.
Consequently, the split real form of $G_2$ does not admit a Shimura datum.
\end{theorem}

\begin{proof}
A connected reductive group $G$ over $\Q$ admits a Shimura datum
if its symmetric space $G(\R)^0 / K$ is a Hermitian symmetric domain
(Deligne, 1979).
The irreducible Hermitian symmetric domains correspond to groups of types
AIII, BDI (rank~2), CI, DIII, EIII, and EVII; no real form of type $G_2$ appears
in this classification (see Helgason, \emph{Differential Geometry, Lie Groups,
and Symmetric Spaces}, Table~V, Chapter~X).
For the split real form $G_2^{\mathrm{split}}$, the maximal compact subgroup is
$K \cong \SO(4)$ and $\dim(G_2/K) = 14 - 6 = 8$.
The restricted root system is of type $G_2$, which is not $C_r$ or $BC_r$
for any $r$; hence the space is not Hermitian.
This is formalized in \texttt{Bridge.lean}, theorem \texttt{no\_shimura\_at\_octonionic\_level}.
\end{proof}

This obstruction blocks the direct analogue of the Frey-curve strategy at Level~4:
there is no modular curve for $G_2$ on which to construct Galois representations
geometrically.

\subsection{The Gross-Savin program}

The viable route: the Gross-Savin program~\cite{GrossSavin1998}.
Automorphic forms on compact $G_2$ theta-lift to Siegel cusp forms on $\GSp_6$,
which \emph{is} a Shimura variety.
Galois representations then exist via Kret--Shin~\cite{KretShin2023},
with Satake parameters forcing the image into $G_2(\Q_\ell) \subset \GSp_6(\Q_\ell)$.

Current status:
\begin{itemize}
  \item Steps 1--2 (theta lift to nonzero Siegel cusp form): proved by
    Gross--Savin~\cite{GrossSavin1998}.
  \item Step 3 (Galois representation construction): proved by
    Kret--Shin~\cite{KretShin2023}, from the general $\mathrm{GSpin}$ construction.
  \item Step 4 (motivic Galois group is $G_2$): conditional, by
    Cauchi--Lemma--Rodrigues Jacinto~\cite{CLRJ2025},
    pending non-vanishing of an archimedean integral.
\end{itemize}
Global modularity lifting ($R = T$ for $G_2$) remains wide open.

\subsection{Connection to Bhargava's higher composition laws}

Split $G_2$ acts on $\Z^2 \otimes \Z^3 \otimes \Z^3$ (Zorn vector matrices),
parametrizing cubic rings.
Bhargava~\cite{Bhargava2004} showed $\GL_2(\Z) \times \GL_3(\Z)$-orbits
on $2 \times 3 \times 3$ integer cubes parametrize pairs of ideal classes in cubic rings.

The octonionic Fermat solutions live on the trilinear variety
$V(\varphi) = \{ (a, b, c) \in (\Z^7)^3 : \varphi(a, b, c) = 0 \}$
where $\varphi(a, b, c) = \langle c, a \times b \rangle$ is the $G_2$-invariant
trilinear form.
Computational pipeline: 6{,}560 Zorn matrices $\to$ 9{,}565 cubic forms
$\to$ 47{,}961 cubic rings, all with discriminant $\equiv 0 \pmod{4}$.

\subsection{The magic square}

The Freudenthal-Tits magic square~\cite{FreudenthalTits} maps pairs of
division algebras to Lie algebras via
$M(A, B) = \mathrm{Der}(A) \oplus (A_0 \otimes J_0) \oplus \mathrm{Der}(J)$
($J = J_3(B)$):
\[
\begin{array}{c|cccc}
  & \R & \C & \HH & \OO \\ \hline
  \R & A_1 & A_2 & C_3 & F_4 \\
  \C & A_2 & A_2 \oplus A_2 & A_5 & E_6 \\
  \HH & C_3 & A_5 & D_6 & E_7 \\
  \OO & F_4 & E_6 & E_7 & E_8
\end{array}
\]
The octonionic column produces all five exceptional Lie algebras.
$M(\OO, \OO) = E_8$ connects to the $E_8$ lattice
(Viazovska~\cite{Viazovska2017}) and, speculatively, to exceptional Langlands.

%% ═══════════════════════════════════════════════════════════════════════
%% 10. MACHINE VERIFICATION
%% ═══════════════════════════════════════════════════════════════════════

\section{Machine verification}\label{sec:lean}

Every theorem in this paper is formalized in Lean~4 with Mathlib.

\begin{table}[h]
\centering
\caption{Key files in the Lean formalization}\label{tab:files}
\begin{tabular}{@{}llc@{}}
  \toprule
  File & Content & Key Theorems \\
  \midrule
  \texttt{Algebra.lean} & Octonion/quaternion integer definitions & --- \\
  \texttt{OctonionicFermat.lean} & Cubic verification (\texttt{native\_decide}) & Thm.~\ref{thm:cubic} \\
  \texttt{SymbolicFermatGeneral.lean} & Power reduction (kernel-checked) & Thm.~\ref{thm:power-reduction} \\
  \texttt{AllOddPrimes.lean} & All odd $p$ (kernel-checked) & Thm.~\ref{thm:allodd} \\
  \texttt{DimensionThreshold.lean} & Dimension boundary & Thm.~\ref{thm:dimension} \\
  \texttt{ZeroDivisorFermat.lean} & Zero-divisor collapse (kernel-checked) & Thm.~\ref{thm:zerodiv} \\
  \texttt{SedenionZeroDivisor.lean} & Sedenion obstruction analysis & Remark~7.3 \\
  \texttt{G2OrbitCollapse.lean} & $G_2$ orbit analysis & Thm.~\ref{thm:g2orbit} \\
  \texttt{Bridge.lean} & Langlands stratification & Thm.~\ref{thm:shimura} \\
  \texttt{Conjectures.lean} & 17 conjectures (documentation) & --- \\
  \bottomrule
\end{tabular}
\end{table}

\noindent
Totals: 29 files, 9{,}102 lines, 355 declarations (theorems, lemmas, and definitions).
The formalization is available at
\url{https://github.com/awktavian/fermat} (branch \texttt{main},
Lean toolchain \texttt{v4.16.0}, Mathlib \texttt{2026-03-20}).

The formalization uses two proof strategies:
\begin{enumerate}
  \item \textbf{Kernel-checked symbolic proofs} for the main theorems
    (Theorems~\ref{thm:power-reduction}, \ref{thm:allodd}, \ref{thm:zerodiv}).
    These use only ring axioms and structural induction, with no \texttt{native\_decide}
    and no \texttt{sorry}.
    They are valid in any ring satisfying the hypotheses.
  \item \textbf{Computational verification} (\texttt{native\_decide}) for concrete instances:
    specific triples in $\HH_\Z$ and $\OO_\Z$, all 21 family verifications,
    and the Gaussian integer exhaustive search.
    These verify that the abstract hypotheses are satisfied by the concrete algebras.
    The sedenion analysis (Section~\ref{sec:sedenion}) is partially formalized;
    the obstruction $\alpha^2 = -2$ was identified during formalization.
\end{enumerate}

Zero \texttt{sorry}. Zero axioms beyond Lean's foundations
(\texttt{propext}, \texttt{Quot.sound}, \texttt{Classical.choice}).

%% ═══════════════════════════════════════════════════════════════════════
%% 11. OPEN PROBLEMS AND CONJECTURES
%% ═══════════════════════════════════════════════════════════════════════

\section{Open problems and conjectures}\label{sec:open}

Ten conjectures, selected from 17 in \texttt{Conjectures.lean}.
All are precise and falsifiable.

\begin{conjecture}[Artin Confinement --- A4]\label{conj:artin}
Every solution to $a^p + b^p = c^p$ in $\OO_\Z$ ($p$ odd, $\geq 3$) lies
in an associative subalgebra $\cong \HH$.
\end{conjecture}

\noindent
\emph{Evidence.} Over 600{,}000 solutions checked; zero have nonzero associator.
The orbit analysis found 61{,}488 solutions at $\Norm{} \leq 8$, all with
$[a, b, c] = 0$.
\emph{Consistent with} Artin's theorem, which guarantees that any \emph{two}
elements of an alternative algebra generate an associative subalgebra.
The conjecture asserts a stronger \emph{triplewise} confinement:
the entire solution triple $(a, b, c)$ lies in a single quaternionic subalgebra.

\begin{conjecture}[Waring Tightness --- B1]\label{conj:waring}
$g(3, \HH_\Z) = 4$.
\end{conjecture}

\noindent
\emph{Evidence.} Verified for all 4{,}785 Lipschitz integers with $\Norm{q} \leq 30$.
Exactly 12 elements (all at norm~26) require 4 cubes.

\begin{conjecture}[Split $G_2$ Parametrization --- D1]\label{conj:splitg2}
Split $G_2(\Z)$-orbits on Zorn integer matrices parametrize cubic rings
with discriminant $\equiv 0 \pmod{4}$, via the embedding
$\mathrm{Zorn}(\Z) \hookrightarrow \Z^2 \otimes \Z^3 \otimes \Z^3$.
\end{conjecture}

\noindent
\emph{Evidence.} 6{,}560 Zorn matrices yield 9{,}565 cubic forms and 47{,}961 cubic rings,
all with discriminant $\equiv 0 \pmod{4}$, all associative.

\begin{conjecture}[$G_2$ Fourier Coefficients and Fermat Families --- E4]\label{conj:fourier}
The theta lift of automorphic forms on $G_2^c$ to $\GSp_6$ produces
Siegel modular forms whose Fourier coefficients encode the number of
octonionic Fermat solution families at each norm level.
\end{conjecture}

\begin{conjecture}[$G_2$ Selmer Averaging --- E2]\label{conj:selmer}
$\mathrm{avg}\,|\Sel_2(G_2, V)| = 3$, by analogy with the Bhargava-Shankar theorem
for elliptic curves.
\end{conjecture}

\noindent
\emph{Status.} Speculative. $G_2$ Selmer groups have not been explicitly defined
in the literature; our \texttt{G2SelmerGroup.lean} is a first attempt.

\begin{conjecture}[Hasse Principle for $V(\varphi)$ --- C2]\label{conj:hasse}
For the primitive subfamily $V(\varphi)_{\mathrm{prim}}$
(where $a \times b \neq 0$ and $a, b, c$ are pairwise non-proportional),
the Hasse principle holds: $V(\varphi)_{\mathrm{prim}}(\Q) \neq \varnothing$
if and only if $V(\varphi)_{\mathrm{prim}}(\Q_v) \neq \varnothing$ for all places $v$.
\end{conjecture}

\noindent
\emph{Evidence.} No local-global failures found; the octonionic cross product
is surjective modulo every prime tested ($p \leq 13$),
so no Brauer-Manin obstruction at finite primes.

\begin{conjecture}[Sedenion Fermat --- S3]\label{conj:sedenion}
Does the Fermat equation $a^n + b^n = c^n$ have nontrivial solutions in $\Se_\Z$
for $n \geq 3$?  Mechanism~A gives odd-exponent solutions (via anti-commuting pairs
in quaternionic subalgebras). Mechanism~B fails (Section~\ref{sec:sedenion}).
The even-exponent case is open.
\end{conjecture}

\begin{conjecture}[Waring Linear Growth --- B2]\label{conj:waringlinear}
$g(p, \HH_\Z) \sim p + 1$ for odd primes $p$.
\end{conjecture}

\noindent
\emph{Evidence.} Three data points: $g(3) \leq 4$, $g(5) \leq 5$, $g(7) \leq 6$.

\begin{conjecture}[Zariski Density --- C3]\label{conj:zariski}
$V(\varphi)(\Q)$ is Zariski-dense in $V(\varphi)$.
\end{conjecture}

\noindent
\emph{Argument.} The fiber structure $\varphi(a, b, c) = \langle c, a \times b \rangle$
gives codimension-1 linear fibers in $c$ for each $(a, b)$ with $a \times b \neq 0$.
Dense base and dense fibers imply dense total space.

\begin{conjecture}[$\varphi$-Contraction as $G_2$ Curvature --- F3]\label{conj:curvature}
The symmetrized contraction identity for the octonionic trilinear form $\varphi$
(proved algebraically in \texttt{TrilinearVariety.lean})
encodes the curvature of $G_2$-holonomy manifolds in the sense of
Bryant~\cite{Bryant1987} and Joyce~\cite{Joyce2000}.
\end{conjecture}

%% ═══════════════════════════════════════════════════════════════════════
%% REFERENCES
%% ═══════════════════════════════════════════════════════════════════════

\begin{thebibliography}{99}

\bibitem{Baez2002}
J.~C.~Baez,
\emph{The octonions},
Bull.\ Amer.\ Math.\ Soc.\ \textbf{39} (2002), 145--205.

\bibitem{Bhargava2004}
M.~Bhargava,
\emph{Higher composition laws I: A new view on Gauss composition,
and quadratic generalizations},
Ann.\ Math.\ \textbf{159} (2004), 217--250.

\bibitem{BhargavaShankar2015}
M.~Bhargava and A.~Shankar,
\emph{Binary quartic forms having bounded invariants, and the boundedness
of the average rank of elliptic curves},
Ann.\ Math.\ \textbf{181} (2015), 191--242.

\bibitem{Bryant1987}
R.~L.~Bryant,
\emph{Metrics with exceptional holonomy},
Ann.\ Math.\ \textbf{126} (1987), 525--576.

\bibitem{BZSV2024}
A.~Ben-Zvi, Y.~Sakellaridis, and A.~Venkatesh,
\emph{Relative Langlands duality},
preprint (2024).

\bibitem{CLRJ2025}
A.~Cauchi, F.~Lemma, and J.~Rodrigues Jacinto,
\emph{Algebraic cycles and functorial lifts from $G_2$ to $\mathrm{PGSp}_6$},
Algebra Number Theory (2025).

\bibitem{ConwaySmith2003}
J.~H.~Conway and D.~A.~Smith,
\emph{On Quaternions and Octonions},
A.~K.~Peters, 2003.

\bibitem{ElkiesGross2001}
N.~Elkies and B.~Gross,
\emph{Cubic rings and the exceptional Jordan algebra},
Duke Math.\ J.\ \textbf{109} (2001), 383--409.

\bibitem{FarguesScholze2021}
L.~Fargues and P.~Scholze,
\emph{Geometrization of the local Langlands correspondence},
arXiv:2102.13459 (2021).

\bibitem{FreudenthalTits}
H.~Freudenthal, \emph{Beziehungen der $E_7$ und $E_8$ zur Oktavenebene, I--XI},
Indagationes Math.\ (1954--1964);
J.~Tits,
\emph{Alg\`ebres alternatives, alg\`ebres de Jordan et alg\`ebres de Lie exceptionnelles},
Indag.\ Math.\ \textbf{28} (1966), 223--237.

\bibitem{Gaitsgory2024}
D.~Gaitsgory et al.,
\emph{Proof of the geometric Langlands conjecture} (five-part series),
arXiv:2405.03599 et seq.\ (2024).

\bibitem{GanGrossSavin2002}
W.~T.~Gan, B.~Gross, and G.~Savin,
\emph{Fourier coefficients of modular forms on $G_2$},
Duke Math.\ J.\ \textbf{115} (2002), 105--169.

\bibitem{GrossSavin1998}
B.~Gross and G.~Savin,
\emph{Motives with Galois group of type $G_2$: an exceptional theta-correspondence},
Compositio Math.\ \textbf{114} (1998), 153--217.

\bibitem{HarveyLawson1982}
R.~Harvey and H.~B.~Lawson,
\emph{Calibrated geometries},
Acta Math.\ \textbf{148} (1982), 47--157.

\bibitem{Hurwitz1898}
A.~Hurwitz,
\emph{\"Uber die Composition der quadratischen Formen von beliebig vielen Variablen},
Nachr.\ Ges.\ Wiss.\ G\"ottingen (1898), 309--316.

\bibitem{Joyce2000}
D.~Joyce,
\emph{Compact Manifolds with Special Holonomy},
Oxford University Press, 2000.

\bibitem{KretShin2023}
A.~Kret and S.~W.~Shin,
\emph{Galois representations for general symplectic groups},
J.~Eur.\ Math.\ Soc.\ (2023).

\bibitem{Krutelevich2007}
S.~Krutelevich,
\emph{Jordan algebras, exceptional groups, and Bhargava composition},
J.~Algebra \textbf{314} (2007), 924--977.

\bibitem{LesliePollack2024}
S.~Leslie and A.~Pollack,
\emph{Modular forms of half-integral weight on exceptional groups},
Compositio Math.\ \textbf{160} (2024).

\bibitem{LiYuan2025}
H.~Li and P.~Yuan,
\emph{Fermat's and Catalan's equations over $M_2(\Z)$},
arXiv:2503.03621, preprint (2025).

\bibitem{Moreno1998}
G.~Moreno,
\emph{The zero divisors of the Cayley-Dickson algebras over the real numbers},
Bol.\ Soc.\ Mat.\ Mexicana \textbf{4} (1998), 13--28.

\bibitem{APollack2020}
A.~Pollack,
\emph{Modular forms on $G_2$ and their standard $L$-function},
Duke Math.\ J.\ \textbf{169} (2020), 1537--1602.

\bibitem{PPollack2018}
P.~Pollack,
\emph{Waring's problem for integral quaternions},
Indag.\ Math.\ \textbf{29} (2018), 1259--1269.

\bibitem{Reggiani2024}
S.~Reggiani,
\emph{The geometry of sedenion zero divisors},
arXiv:2411.18881, preprint (2024).

\bibitem{Ribenboim1979}
P.~Ribenboim,
\emph{13 Lectures on Fermat's Last Theorem},
Springer, 1979.

\bibitem{Ribenboim2004}
P.~Ribenboim,
\emph{Fermat's equation for matrices or quaternions over $q$-adic fields},
Acta Arith.\ \textbf{113} (2004), 241--250.

\bibitem{Viazovska2017}
M.~Viazovska,
\emph{The sphere packing problem in dimension 8},
Ann.\ Math.\ \textbf{185} (2017), 991--1015.

\bibitem{Wiles1995}
A.~Wiles,
\emph{Modular elliptic curves and Fermat's Last Theorem},
Ann.\ Math.\ \textbf{141} (1995), 443--551.

\end{thebibliography}

\end{document}
